\sgaStableTarget{sga6.e14.expose}\section*{\underline{Exposé XIV}\\[3pt] \underline{\underline{OPEN PROBLEMS IN INTERSECTION THEORY}}}
\begin{center}
\underline{by A. Grothendieck}
\end{center}

In this exposé we review some of the problems connected with the questions encountered in this seminar that deserve to be studied. We leave aside the problems connected with Serre's \uline{local theory} of intersections [14], which we have not had occasion to use in the Seminar, and refer in particular to loc. cit. for the two fundamental conjectures that remain open in this direction (loc. cit. Chap. V \S~4). We also pass over the problems, very important for arithmetic applications, connected with the \uline{structure of groups of cycle classes}, especially on a smooth projective algebraic scheme; for a discussion of these problems we refer to S. KLEIMAN, Algebraic cycles and the Weil conjectures in Dix exposés sur la cohomologie des schémas (Masson-North Holland Publishing Cie). In no.~4 of this exposé we also give some supplements to the Seminar, concerning in particular the \uline{Chow ring}, which had not found a place in the preceding exposés.

\sgaStableTarget{sga6.e14.sec.1}\section*{1. Operations @@M0@@ in the derived category @@M1@@\SGAAuthorFootMark{*}}
\SGAAuthorFootText{*}{cf. the note on page 23.}

Let @@M2@@ be a locally ringed topos, and let @@M3@@ be the ring of classes of perfect complexes on @@M4@@ (Exp. IV). We would like to define on @@M5@@ a @@M6@@-ring structure (Exp. V), functorial in @@M7@@, such that the natural homomorphism
@@M8@@
(where the first member is the ring of classes of locally free Modules on @@M9@@, endowed with the @@M10@@-structure defined in Exp. V) be a @@M11@@-homomorphism. For this purpose, one would need to define functors (generally nonadditive) of ``exterior power''
@@M12@@
taking pseudo-coherent (resp. perfect) complexes to pseudo-coherent (resp. perfect) complexes and having the formal properties needed, with respect to exact triangles, to define a @@M13@@-structure on @@M14@@ by the formula
\sgaStableTarget{sga6.e14.eq.1-1}@@M15@@

